\section{Parte 2}

\begin{ejercicio} % // TODO: Ejercicio 5 del Final 2025
    Para $n\geq 10^6$, se tiene $x_n:\left]-1,\nicefrac{1}{2}\right[\to \mathbb{R}$ que cumple:
    \begin{equation*}
        x_n(t) = \frac{1}{n} + \int_{0}^{t} e^{x_n(s)}~ds , \qquad t\in \left]-1,\frac{1}{2}\right[
    \end{equation*}
    Demostrar que $\{x_n\}$ converge uniformemente en $\left]-1,\frac{1}{2}\right[$.\\

    \noindent
    Pasamos la ecuación integral a un problema de valores iniciales:
    \begin{equation*}
        \dot{x} = e^x, \qquad x(0) = \frac{1}{n}
    \end{equation*}
    Tenemos unicidad para toda condición inicial. Resolvemos el PVI para la condición $x(0) = 0$, por variables separadas:
    \begin{equation*}
        -e^{-x} = \int e^{-x}~dx = \int dt = t + c \quad\Longrightarrow\quad e^{-x} = 1-t \quad\Longrightarrow\quad x = -\ln(1-t), \quad t<1
    \end{equation*}
    Aplicando la condición inicial tenemos $c = -1$. Consideramos así la solución:
    \begin{equation*}
        x(t) = -\ln(1-t), \qquad t\in [-1,\nicefrac{1}{2}]
    \end{equation*}
    Aplicando el Teorema de la dependencia continua respecto a condiciones iniciales, existe $N\in \mathbb{N}$ de manera que si $n\geq N$ entonces existe $x_n(t)$ solución del PVI para la condición inicial $x(0) = \frac{1}{n}$ y de forma que $\{x_n\}\to x$ uniformemente en $[-1,\frac{1}{2}]$, luego también en $\left]-1,\frac{1}{2}\right[$.
\end{ejercicio}

\begin{ejercicio} % // TODO: Jorge dice que esto es un ejercicio propuesto en clase
    Sea $f:\left]a,b\right[\times [c,d]\to \mathbb{R}$, $f=f(t,\lm)$. Son equivalentes:
    \begin{enumerate}
        \item[$i)$] $f(t,\lm)\to 0$ converge uniformemente si $t\to a^+$ en $\lm \in [c,d]$.
        \item[$ii)$]$\forall t_n\to a^-$, $f(t_n,\lm)\to 0$ converge uniformemente en $\lm \in [c,d]$.
    \end{enumerate}
    Veámoslo:
    \begin{description}
        \item [$(i)\Longrightarrow  ii)$] Por definición del límite uniforme:
            \begin{equation*}
                \forall \varepsilon>0~~\exists \delta>0~:~ 0 < t-a<\delta \Longrightarrow  |f(t,\lm)| < \varepsilon\quad \forall \lm \in [c,d]
            \end{equation*}
            Así, tomando $t_n\to a^-$, existe $N\in \mathbb{N}$ de forma que si $n\geq N$ entonces:
            \begin{equation*}
                t_n-a < \delta \quad\Longrightarrow\quad |f(t_n,\lm)| < \varepsilon \quad \forall \lm \in [c,d]
            \end{equation*}
        \item [$(ii)\Longrightarrow (i)$]  Por contrarrecíproco, suponemos que no hay convergencia uniforme, es decir, que existe $\varepsilon>0$ tal que $\forall \delta>0$ existen $t$ con $0<t-a<\delta$ y $\lm \in [c,d]$ tales que:
            \begin{equation*}
                |f(t,\lm)| \geq \varepsilon
            \end{equation*}
            Ahora, para cada $n\in \mathbb{N}$ tomamos $\delta = \frac{1}{n}$, con lo que existen $t_n$ con $0<t_n-a<\delta = \frac{1}{n}$ y $\lm_n\in [c,d]$ tales que:
            \begin{equation*}
                |f(t_n,\lm_n)| \geq \varepsilon
            \end{equation*}
            Y como $t_n\to a^-$ tenemos que:
            \begin{equation*}
                \sup_{\lm} |f(t_0,\lm)| \geq |f(t_n,\lm_n)| \geq \varepsilon
            \end{equation*}
    \end{description}
\end{ejercicio}

\begin{ejercicio} % // TODO: También un ejercicio de clase
    Tenemos $f_n:[a,b]\to \mathbb{R}$ continuas con $\{f_n\}\to f$ uniformemente y tenemos $\{t_n\}\to t_\ast$ con $t_n \in [a,b]$, entonces $\{f_n(t_n)\}\to f(t_\ast)$.\\

    \noindent
    En primer lugar tenemos por la convergencia uniforme que $f:[a,b]\to \mathbb{R}$  es continua. Ahora, para $t_\ast \in [a,b]$ tenemos:
    \begin{equation*}
        |f_n(t_n)-f(t_\ast)| \leq |f_n(t_n)-f(t_n)| + |f(t_n)-f(t_\ast)| \leq \|f_n-f\|_\infty + |f(t_n)-f(t_\ast)|
    \end{equation*}
    Con cada sumando convergiendo a cero:
    \begin{equation*}
        \{f_n\}\to f \text{\ uniformemente} \quad\Longleftrightarrow\quad \sup_{t\in [a,b]}|f_n(t)-f(t)| = \|f_n-f\|_\infty \to 0
    \end{equation*}
    El primero por la convergencia uniforme y el segundo por la continuidad de $f$.

    % // TODO: Ortega deja planteado el ejercicio recíproco:
    % Cómo ir de {f_n(t_n)} --> f(t_*) => {f_n} --> f uniformemente
\end{ejercicio}

\begin{ejercicio} % // TODO: Ejercicio 5 del parcial 2  del 24
    Se considera $U\in C^1(\mathbb{R}^2)$ y el sistema:
    \begin{equation*}
        \left\{\begin{array}{l}
                \ddot{x} = \frac{\partial U}{\partial x}(x,y) \\
                \\
                \ddot{y} = \frac{\partial U}{\partial y}(x,y) 
        \end{array}\right.
    \end{equation*}
    \begin{enumerate}[label=\alph*)]
        \item Dada una solución $(x(t),y(t))$, demostrar que $\exists a,b,c\in \mathbb{R}\setminus \{0\}$ de forma que:
            \begin{equation*}
                E(t) = a\dot{x}(t)^2 + by(t)^2 + cu(x(t),y(t))
            \end{equation*}
            es constante.\\

            \noindent
            Para ello, calculamos primero la derivada de $E$, igualamos a cero y despejamos, obteniendo que para $\lm\in \mathbb{R}\setminus \{0\}$ se debe tener:
            \begin{equation*}
                a = \lm = b, \qquad c = -2\lm
            \end{equation*}
        \item Supuesto que $U(x,y)\leq -3\quad \forall (x,y)\in \mathbb{R}^2$, probar que toda solución maximal está definida hasta $+\infty$.

            Buscamos aplicar el Teorema de Prolongación. Para ello, el primer paso es pasar el sistema anterior a primer orden, tomando:
            \begin{equation*}
                z = (z_1,z_2,z_3,z_4) \in \mathbb{R}^4, \qquad z = Z(t,z)
            \end{equation*}
            donde:
            \begin{equation*}
                Z(t,z) = \left(z_3,z_4,\frac{\partial U}{\partial z_1}(z_1,z_2),\frac{\partial U}{\partial z_2}(z_1,z_2)\right)
            \end{equation*}
            Es un campo que está bien definido y es continuo en $D=\mathbb{R}\times\mathbb{R}^4$. Sea $z(t) = (x(t),y(t),\dot{x}(t),\dot{y}(t))$ solución maximal de $\dot{z} = Z(t,z)$ definido en $\left]\alpha,\omega\right[$. Supuesto que $\omega<\infty$, el Teorema de Prolongación nos dice que:
            \begin{itemize}
                \item Bien $\|z(t)\|\to\infty$ cuando $t\to \omega^-$.
                \item $\exists t_n\to\omega, \xi \in \mathbb{R}^4$ con $z(t_n)\to \xi$, $(\omega,\xi)\in \partial D$.
            \end{itemize}
            La segunda opción se descarta por tener $\partial D =\emptyset $. Para la primera opción, tomando $\lm=1$ tenemos que:
            \begin{equation*}
                \mathbb{R}\ni E_0 = \dot{x}(t)^2 + \dot{y}(t)^2 -2U(x(t),y(t))     
            \end{equation*}
            de aquí se llega a que:
            \begin{equation*}
                \dot{x}(t)^2 + \dot{y}(t)^2 = E_0 + 2U(x(t),y(t)) \leq E_0-6
            \end{equation*}
            Por tanto:
            \begin{equation*}
                \|(\dot{x}(t),\dot{y}(t))\|_2 \leq \sqrt{E_0-6} = M \qquad \forall t\in \left]\alpha,\omega\right[
            \end{equation*}
            Ahora, para $t_0\in \left]\alpha,\omega\right[$ se tiene que:
            \begin{equation*}
                (x(t),y(t)) = (x(t_0),y(t_0)) + \int_{t_0}^{t} (\dot{x}(s), \dot{y}(s))~ds 
            \end{equation*}
            de donde:
            \begin{align*}
                \|(x(t),y(t))\| &\leq \|(x(t_0),y(t_0))\| + \left|\int_{t_0}^{t} (\dot{x}(s), \dot{y}(s))~ds \right| \leq \|x(t_0),y(t_0)\| + M|t-t_0| \\
                                &\leq \|(x(t_0)),y(t_0))\| + M(\omega - t_0) \qquad \forall t\in \left[t_0,\omega\right[
            \end{align*}
            Así, si en $\mathbb{R}^4$ consideramos la norma:
            \begin{equation*}
                \|(x,y,z,t)\|_4 = \|(x,y)\|_2 + \|(z,t)\|_2
            \end{equation*}
            Hemos acabado.
    \end{enumerate}
\end{ejercicio}

% // TODO: Ejercicios de Ortega
\begin{ejercicio}
    Tenemos el sistema:
    \begin{equation*}
        \left\{\begin{array}{ll}
                \dot{x} = y^2, & x(0) = x_0 \\
                \dot{y} = x-1, & y(0) = y_0
        \end{array}\right.
    \end{equation*}
    Es un sistema autónomo y se nos pide calcular:
    \begin{equation*}
        \frac{\partial x}{\partial y_0}(2;1,0)
    \end{equation*}
    Claramente buscamos aplicar el Teorema de diferenciabilidad respecto a condiciones iniciales. Se nos pregunta por la primera coordenada de la solución respecto a la segunda coordenada de la condición inicial. Escribimos:
    \begin{equation*}
        z = \left(\begin{array}{cc}
            x \\
            y 
        \end{array}\right) \in \mathbb{R}^2
    \end{equation*}
    Con lo que tenemos el sistema $\dot{z} = Z(t,z), z(0) = z_0 = \left(\begin{array}{cc}
        x_0 \\
        y_0 
    \end{array}\right)$ donde:
    \begin{equation*}
        Z(t,z) = \left(\begin{array}{cc}
            y^2 \\
            x-1 
        \end{array}\right)
    \end{equation*}
    Observamos que las dos coordenadas son de clase $C^1$, con lo que podemos aplicar el Teorema de diferenciabilidad respecto a condiciones iniciales para el dominio $D=\mathbb{R}\times\mathbb{R}^2$. Sabemos que la solución general $z(t;t_0,z_0)\in C^1(\cc{D},\mathbb{R}^2)$. Si escribimos:
    \begin{equation*}
        \frac{\partial z}{\partial z_0} = \left(\begin{array}{cc}
                \frac{\partial x}{\partial x_0} & \red{\frac{\partial x}{\partial y_0}} \\
            \\
            \frac{\partial y}{\partial x_0} &  \frac{\partial y}{\partial y_0}
        \end{array}\right)
    \end{equation*}
    Por lo que nos interesa la coordenada $(1,2)$ de la matriz. Escribimos ahora:
    \begin{equation*}
        Y(t) = \frac{\partial z}{\partial z_0}(t;0,z_0)
    \end{equation*}
    con lo que obtenemos la ecuación variacional ($t_0 = 0$):
    \begin{equation*}
        \dot{Y} = \frac{\partial Z}{\partial z_0}(t,z(t;t_0,z_0))\ Y, \qquad Y(t_0) = I
    \end{equation*}
    Observamos que el sistema admite el equilibrio (solución constante) $(x\equiv1,y\equiv0)$, es decir:
    \begin{equation}\label{eq:ej_ref}
        x(t;t_0=0,1,0) = 1, \qquad y(t;t_0=0,1,0) = 0, \quad t\in \mathbb{R}
    \end{equation}
    Ahora:
    \begin{equation*}
        \frac{\partial Z}{\partial z} = \left(\begin{array}{cc}
            0 & 2y \\
            1 & 0 
        \end{array}\right)
    \end{equation*}
    Y tenemos que:
    \begin{equation*}
        \dot{Y} = \left(\begin{array}{cc}
            0 & 0 \\
            1 & 0 
        \end{array}\right)\ Y, \qquad Y(0) = I
    \end{equation*}
    Buscamos primero soluciones vectoriales de este sistema, $y = (y_1,y_2)$:
    \begin{equation*}
        \left\{\begin{array}{l}
                \dot{y}_1 = 0 \\
                \dot{y}_2 = y_1
        \end{array}\right.
    \end{equation*}
    por lo que:
    \begin{equation*}
        y_1(t) = c_1, \qquad y_2(t) = c_2 + c_1 t
    \end{equation*}
    imponiendo condiciones iniciales, $c_1 = 1, c_2 = 0$.
    de donde:
    \begin{equation*}
        Y(t) = \left(\begin{array}{cc}
            1 & 0 \\
            t & 1 
        \end{array}\right)
    \end{equation*}
    En definitiva:
    \begin{equation*}
        \frac{\partial x}{\partial y_0}(t;1,0) = 0
    \end{equation*}
    En particular vale 0 para $t=2$, sabemos que $(2;1,0)\in \cc{D}$, pues las soluciones~\eqref{eq:ej_ref} están definidas en todo $\mathbb{R}$.
\end{ejercicio}

\begin{ejercicio}
    Cambiemos el problema anterior, tenemos el sistema:
    \begin{equation*}
        \left\{\begin{array}{ll}
                \dot{x} = y^2, & x(0) = \varepsilon+1 \\
                \dot{y} = x-1, & x(0) = \varepsilon
        \end{array}\right.
    \end{equation*}
    Se pide calcular $\frac{\partial y}{\partial \varepsilon}(2;0)$.
\end{ejercicio}
